\begin{trapbox}
\textbf{\textcolor{CTrap}{原则：每个容易误解的地方，都给出一个能算清楚的小例子。}}


\begin{center}
\begin{tikzpicture}[x=.9cm,y=.9cm,every node/.style={font=\small}]
\draw[blue!75!black,thick] (0,0) circle (2);
\draw[orange!80!black,thick] (-2.2,0)--(2.3,0);
\fill (1,0) circle (2pt) node[below=5pt] {$R(1,0)$};
\fill (-2,0) circle (2pt); \fill (2,0) circle (2pt);
\node at (0,2.7) {圆内点作出的这条线是割线};
\node at (0,-2.8) {穿过圆周两次，不能叫切线};
\begin{scope}[xshift=7cm]
\draw[blue!75!black,thick] (0,0) circle (2);
\draw[orange!80!black,thick] (2,-2.2)--(2,2.2);
\draw[gray!70] (0,0)--(2,0);
\draw (1.8,0)--(1.8,.2)--(2,.2);
\fill (2,0) circle (2pt) node[below left=5pt] {$T(2,0)$};
\node at (0,2.7) {真正切线只接触圆周一次};
\node at (0,-2.8) {半径与切线垂直};
\end{scope}
\end{tikzpicture}

\end{center}
左右使用同一半径与尺度。圆内点所在直线穿过圆周，不能冒充只接触一次的切线；相应极线 $x=4$ 见直觉理解的三态图。

\begin{description}[style=nextline,leftmargin=2.8em,itemsep=0.65em]
\item[\textbf{易错一：把圆内点的极线当成两切点的连线。}]
\textbf{错误想法：} 既然极线是切点弦，任何点都应当能画出两切点。

\textbf{最小反例：} 圆 $x^2+y^2=4$，取圆内点 $R(1,0)$。点 $R$ 到圆心的距离是 $1<2$，从它不能向圆作真实切线；但按统一规则，$l_R:x=4$ 仍是一条直线，且它与圆没有公共点。

\textbf{正确说法：} “圆外\textbf{可以}用真实切点弦解释；圆上\textbf{可以}用真实切线解释；圆内\textbf{只能先用统一方程解释}。”三种情况公式相同，图形含义不完全相同。

\item[\textbf{易错二：把极点和极线当成同一种对象。}]
\textbf{最小反例：} 圆 $x^2+y^2=4$ 中，点 $P(4,0)$ 是\textbf{点}，直线 $x=1$ 是\textbf{线}；它们构成一对对应关系，不能说“极点 $x=1$”或“极线 $(4,0)$”。

\textbf{正确说法：} $P$ 是 $x=1$ 对应的极点；$x=1$ 是 $P$ 的极线。若出现“共轭”，只把它理解成\textbf{两点互在对方的极线上}，不用引入另外一套术语。

\item[\textbf{易错三：把“$Q$ 在 $P$ 的极线上”看成“$Q$ 在圆上”。}]
\textbf{最小反例：} 对圆 $x^2+y^2=4$，$P(4,0)$ 的极线为 $x=1$。点 $Q(1,0)$ 满足 $x=1$，却有 $1^2+0^2=1\ne4$，不在圆上。

\textbf{正确说法：} \textbf{点是否在极线上}，代入极线方程；\textbf{点是否在曲线上}，代入曲线方程。两条方程不是同一个。

\item[\textbf{易错四：交换谁的极线、谁在极线上时，下标弄反。}]
\textbf{错误写法：} “$P$ 在 $Q$ 的极线上，所以直接代 $P$ 到 $P$ 的极线”。这把点和直线的所属关系混淆了。

\textbf{正确检查：} $Q\in l_P$ 应当\textbf{将 $Q$ 的坐标代入 $l_P$}；$P\in l_Q$ 则将 $P$ 的坐标代入 $l_Q$。当两式成立与否相同时，才叫互换关系。对圆，其共同检验式就是 $x_1x_2+y_1y_2=r^2$。

\item[\textbf{易错五：以为随便一条固定直线都能对应到一个普通平面内的点。}]
\textbf{最小反例（圆）：} 对圆 $x^2+y^2=4$，任意点 $P(u,v)$ 的极线方程为 $ux+vy=4$。直线 $x=0$ 经过圆心。若想让 $ux+vy=4$ 与 $x=0$ 是同一条直线，就必须有 $v=0$ 且右边常数为 $0$，但实际是 $4$，不可能。因而这条过圆心的直线\textbf{不是某个普通有限点的极线}。

\textbf{正确用法：} 先确认 $l=l_P$ 确实成立，再从 $Q\in l_P$ 推出 $P\in l_Q$。不能把推论直接套在任何给定的直线上。

\item[\textbf{易错六：忘记圆、椭圆、双曲线中心点的例外。}]
\textbf{最小反例：} 圆 $x^2+y^2=4$，取 $P(0,0)$，会得到 $0x+0y=4$，也就是 $0=4$。这\textbf{不是直线}，所以不能称为圆心的普通极线。椭圆、双曲线中心同样会出现 $0=1$。

\textbf{正确说法：} 对这些标准方程，中心点不能直接参加本节的“一个点对应一条普通直线”的描述。不要说“中心的极线就是坐标轴”或“中心的极线是任意直线”。

\item[\textbf{易错七：把圆的方程误用到椭圆，或者漏掉双曲线的负号。}]
\textbf{最小反例：} 圆 $x^2+y^2=4$，$Q(1,1)$ 的极线是 $x+y=4$；但椭圆 $\frac{x^2}{4}+y^2=1$，同一个 $Q(1,1)$ 的极线却是 $\frac{x}{4}+y=1$。它们\textbf{不同}。双曲线 $\frac{x^2}{A^2}-\frac{y^2}{B^2}=1$ 的极线中，第二项仍应带负号。

\textbf{正确做法：} 先认出曲线类型，再写对应公式；不要仅凭“极线”二字套同一个系数。

\item[\textbf{易错八：把抛物线的 $x+x_1$ 写成 $2xx_1$。}]
\textbf{最小反例：} 抛物线 $y^2=2px$，点 $P(x_1,y_1)$ 的极线是 $y_1y=p(x+x_1)$。因为原方程关于 $x$ 只有\textbf{一次项}，所以要出现的是两个横坐标的\textbf{和}，而不是两个横坐标的积。

\textbf{检查法：} 取 $p=1$、$P(2,2)$，应该得到 $2y=x+2$，将 $Q(-2,0)$ 代入可以成立；错写 $2xx_1$ 就无法保持这一关系。

\item[\textbf{易错九：以为所有动点 $Q$ 的极线都会穿过 $Q$ 本身。}]
\textbf{最小反例：} 圆 $x^2+y^2=4$，$Q(1,0)$ 的极线是 $x=4$，而 $Q$ 的横坐标为 $1$，所以 $Q\notin l_Q$。

\textbf{正确结论：} 对圆，$Q(u,v)\in l_Q\iff u^2+v^2=r^2$，即\textbf{只有 $Q$ 位于圆上时，点才在自己的极线上}。此时对应直线正是该点的切线。椭圆、双曲线、抛物线的标准形式也有“点在曲线上才在自己的极线上”的性质。

\item[\textbf{易错十：只会说“对称”，却不会判断这种方法有没有提速。}]
\textbf{错误习惯：} 见到圆锥曲线就先求极线，结果反而多出无用计算。

\textbf{正确判断：} 题目有“外点引切线--切点弦--动点在定直线--证明直线过定点”等信号，才优先考虑本结论；若只是求普通弦长、交点坐标或者斜率，通常先用课内常规方法。\textbf{本结论是提高拓展，不是所有题的万能公式。}

\item[\textbf{易错十一：一般二次方程中把“一半”丢掉。}]
对 $\alpha x^2+\beta xy+\gamma y^2+\delta x+\varepsilon y+\varphi=0$，写对应直线时，$xy$ 对应 $\frac\beta2(x_1y+y_1x)$，$x$ 对应 $\frac\delta2(x_1+x)$，$y$ 对应 $\frac\varepsilon2(y_1+y)$。\textbf{原因：} 每个一次项都由两个点的坐标平均组成，交换两点后才能保持原式。

\textbf{学习建议：} 普通高中学习阶段优先使用圆、椭圆、双曲线、抛物线各自的简式；这一长式是检查和推广，不必硬背。
\end{description}
\end{trapbox}
